\begin{table*}[t]
\centering
\caption{DeepScript Whole-Inscription Corpus: Epigraphic Provenance, Chronology, and Site-Disjoint Partitions}
\label{tab:dataset_statistics}
\small
\begin{tabular}{l l l l c c c c}
\toprule
\textbf{Script Family} & \textbf{Chronological Period} & \textbf{Geographical Region} & \textbf{Primary Provenance} & \textbf{Train} & \textbf{Val} & \textbf{Test} & \textbf{Total} \\
\midrule
Ashokan Brahmi & c. 3rd C. BCE & Pan-Indian (Pillars & Rocks) & Wikimedia Commons (CC BY-SA / PD) & 6 & 2 & 2 & \textbf{10} \\
Tamil-Brahmi & c. 3rd C. BCE – 5th C. CE & Tamil Nadu & Kerala (Caverns) & TBSI / TNDA (CC BY 4.0) & Wikimedia & 15 & 5 & 5 & \textbf{25} \\
Gupta Script & c. 4th – 8th C. CE & Northern & Central India (Pillars/Slabs) & Wikimedia Commons (CC BY-SA / PD) & 6 & 2 & 2 & \textbf{10} \\
Kharosthi & c. 4th C. BCE – 3rd C. CE & Gandhara / Silk Road (Schist/Reliquaries) & Wikimedia Commons (CC BY-SA / PD) & 7 & 2 & 1 & \textbf{10} \\
Grantha & c. 6th – 14th C. CE & Southern India (Caves & Copper Charters) & Wikimedia Commons (CC BY-SA / PD) & 6 & 2 & 2 & \textbf{10} \\
Kadamba & c. 5th – 6th C. CE & Deccan / Karnataka (Box-Headed Slabs) & Wikimedia Commons (CC BY-SA / PD) & 6 & 2 & 1 & \textbf{9} \\
\midrule
\textbf{Total Inscriptions} & --- & --- & --- & \textbf{46} & \textbf{15} & \textbf{13} & \textbf{74} \\
\bottomrule
\end{tabular}
\vspace{1mm}
\begin{minipage}{\linewidth}
\footnotesize \textit{Note:} Partitioning enforces strict archaeological site-disjointness; no two photographic perspectives of the same rock edict or pillar charter appear across different splits.
\end{minipage}
\end{table*}
